\documentclass[11pt]{article}

\usepackage{amsmath, amssymb}
\usepackage{geometry}
\usepackage{setspace}
\usepackage{tcolorbox} % For attractive definition boxes
\usepackage{booktabs}  % For professional tables
\usepackage{enumitem}  % For better list formatting
\usepackage{xcolor}    % For subtle color styling

\geometry{margin=1in}
\setstretch{1.15}

% Custom styling for headers
\usepackage{titlesec}
\titleformat{\section}{\Large\bfseries\color{blue!80!black}}{\thesection}{1em}{}[\titlerule]
\titleformat{\subsection}{\large\bfseries\color{black!80}}{\thesubsection}{1em}{}

\begin{document}

\begin{center}
\hrule \vspace{0.4cm}
{\Large \textbf{Lecture Scribe: Marginal PMF, Correlation, Covariance, Joint Gaussian RVs}}\\[0.2cm]
\textbf{Course:} CSE400 – Fundamentals of Probability in Computing\\[0.1cm]
\textbf{Lecture:} 18 \textbar\ \textbf{Instructor:} Dhaval Patel, PhD\\[0.1cm]
\textbf{Date:} March 12, 2026\\[0.1cm]
\textbf{Name:} [Student Name] \textbar\ \textbf{Enrollment No:} [Student ID]
\vspace{0.3cm} \hrule
\end{center}

\vspace{0.5cm}

\textbf{\Large Table of Contents}
\begin{enumerate}[label=\textbf{\arabic*.}, leftmargin=*]
    \item Marginal PMF
    \begin{itemize}[label=--]
        \item Definition
        \item Marginal PMF from Joint PMF Table
        \item Solved Example
    \end{itemize}
    \item Correlation
    \item Covariance
    \begin{itemize}[label=--]
        \item Definition and Interpretation
        \item Algebraic Form and Properties
        \item Solved Examples
    \end{itemize}
    \item Pearson’s Correlation Coefficient
    \item Joint Gaussian Random Variables
    \begin{itemize}[label=--]
        \item Definition and Intuition
        \item Algebraic Form and Interpretation
    \end{itemize}
\end{enumerate}

\newpage

\section*{Topic 1: Marginal PMF}

\subsection*{1. Definitions and Notation}

\begin{tcolorbox}[colback=blue!5!white,colframe=blue!75!black,title=\textbf{Definition 1.1 (Joint PMF):}]
For discrete random variables $X$ and $Y$, the joint PMF is defined as:
\[
p_{XY}(x,y) = P(X=x, Y=y)
\]
\end{tcolorbox}

\begin{tcolorbox}[colback=blue!5!white,colframe=blue!75!black,title=\textbf{Definition 1.2 (Marginal PMF):}]
The marginal PMF of a random variable is obtained by summing the joint PMF over all possible values of the other variable.
\[
p_X(x) = \sum_y p_{XY}(x,y), \quad p_Y(y) = \sum_x p_{XY}(x,y)
\]
\end{tcolorbox}

\noindent\textbf{Key Idea:}\\
The marginal PMF of one variable is obtained by summing the joint PMF over all possible values of the other variable.

\subsection*{2. Marginal PMF from Joint PMF Table}

From the table representation on page 7–8:
\begin{itemize}
    \item Row sums $\rightarrow p_X(x)$
    \item Column sums $\rightarrow p_Y(y)$
\end{itemize}

\noindent Mathematically:
\[
p_X(x_i) = \sum_j p_{ij}, \quad p_Y(y_j) = \sum_i p_{ij}
\]

\noindent Also:
\[
\sum_i p_X(x_i) = 1, \quad \sum_j p_Y(y_j) = 1
\]

\noindent\textbf{Explanation:}\\
Summation is applied because probabilities over all possible outcomes must add to 1 (axiom of probability).

\newpage

\subsection*{3. Worked Example 1 (From Lecture PDF)}

\noindent\textbf{Problem Statement (Page 9–11):}\\
Joint PMF is given as:

\begin{table}[h]
\centering
\begin{tabular}{c|cc}
\toprule
\textbf{$p_{XY}(x,y)$} & \textbf{$y=0$} & \textbf{$y=1$} \\
\midrule
$x=0$ & $1/8$ & $1/4$ \\
$x=1$ & $3/8$ & $1/4$ \\
\bottomrule
\end{tabular}
\end{table}

\noindent\textbf{Step-by-Step Solution}

\noindent \textbf{Step 1: Compute $p_X(x)$}
\begin{align*}
p_X(0) &= \frac{1}{8} + \frac{1}{4} = \frac{1}{8} + \frac{2}{8} = \frac{3}{8} \\[0.2cm]
p_X(1) &= \frac{3}{8} + \frac{1}{4} = \frac{3}{8} + \frac{2}{8} = \frac{5}{8}
\end{align*}

\noindent \textbf{Step 2: Compute $p_Y(y)$}
\begin{align*}
p_Y(0) &= \frac{1}{8} + \frac{3}{8} = \frac{4}{8} = \frac{1}{2} \\[0.2cm]
p_Y(1) &= \frac{1}{4} + \frac{1}{4} = \frac{1}{2}
\end{align*}

\noindent \textbf{Final Answer:}
\[
p_X(x) =
\begin{cases}
3/8, & x=0 \\
5/8, & x=1
\end{cases}
\quad
p_Y(y) =
\begin{cases}
1/2, & y=0 \\
1/2, & y=1
\end{cases}
\]

\vspace{0.5cm}
\noindent\textbf{Practice Example (From Sheldon Ross – A First Course in Probability)}

\noindent\textbf{Problem:}\\
Given:
\[
p_{XY}(x,y)=\frac{1}{6} \text{ for } (x,y)\in\{(1,1),(1,2),(2,1),(2,2),(3,1),(3,2)\}
\]
Find marginal PMFs.

\noindent\textbf{Solution:}
\begin{align*}
p_X(1) &= \frac{1}{6}+\frac{1}{6}=\frac{1}{3}, \quad
p_X(2)=\frac{1}{3}, \quad
p_X(3)=\frac{1}{3} \\[0.2cm]
p_Y(1) &= \frac{3}{6}=\frac{1}{2}, \quad
p_Y(2)=\frac{3}{6}=\frac{1}{2}
\end{align*}

\noindent\textbf{Reasoning:} Uniform distribution over 6 outcomes.

\newpage

\section*{Topic 2: Correlation}

\subsection*{1. Definition}

\noindent Correlation is defined as:
\[
\boxed{R_{XY} = E[XY]}
\]

\noindent\textbf{Interpretation (Page 12–13):}
\begin{itemize}
    \item Measures average product of $X$ and $Y$
    \item Computed about origin (not mean)
    \item Depends on scale
    \item Mixes mean, spread, dependence
\end{itemize}

\noindent\textbf{Key Insight:}
\[
R_{XY} = 0 \Rightarrow \text{orthogonal about origin}
\]

\noindent\textbf{Transition:}\\
To remove effect of means $\rightarrow$ use Covariance.

\newpage

\section*{Topic 3: Covariance}

\subsection*{1. Definition and Interpretation}

\[
\text{Cov}(X,Y)=E[(X-\mu_X)(Y-\mu_Y)]
\]

\noindent Where:
\[
\mu_X = E[X], \quad \mu_Y = E[Y]
\]

\noindent\textbf{Interpretation (Page 15–16 diagrams):}
\begin{itemize}
    \item $>0$: Positive covariance $\rightarrow$ variables move together
    \item $<0$: Negative covariance $\rightarrow$ opposite movement
    \item $=0$: No linear relationship
\end{itemize}

\noindent\textbf{Important Note:}\\
Magnitude depends on units $\rightarrow$ cannot measure strength reliably.

\subsection*{2. Algebraic Derivation}

\begin{align*}
\text{Cov}(X,Y) &= E[(X-\mu_X)(Y-\mu_Y)] \\
&= E[XY - X\mu_Y - Y\mu_X + \mu_X\mu_Y] \\
&= E[XY] - \mu_Y E[X] - \mu_X E[Y] + \mu_X\mu_Y \\
&= E[XY] - \mu_X\mu_Y
\end{align*}

\noindent\textbf{Final Formula}
\begin{tcolorbox}[colback=green!5!white,colframe=green!60!black,width=\linewidth,arc=1mm,auto outer arc]
\[
\text{Cov}(X,Y)=E[XY]-E[X]E[Y]
\]
\end{tcolorbox}

\subsection*{3. Properties}

\begin{itemize}
    \item $\text{Cov}(X,Y)=\text{Cov}(Y,X)$
    \item $\text{Cov}(X,X)=\text{Var}(X)$
    \item $\text{Cov}(aX+b,cY+d)=ac\,\text{Cov}(X,Y)$
    \item $X \perp Y \Rightarrow \text{Cov}(X,Y)=0$
\end{itemize}

\noindent\textbf{Important:}\\
Zero covariance $\neq$ independence.

\newpage

\subsection*{4. Worked Example 2 (From Lecture PDF)}

\noindent\textbf{Problem (Page 24–28):}\\
If:
\[
Y=-2X+3
\]
Find:
\[
\text{Cov}(X,Y)
\]

\noindent\textbf{Step-by-step Solution}

\noindent Using:
\[
\text{Cov}(X,Y)=E[XY]-E[X]E[Y]
\]

\noindent \textbf{Step 1: Compute $XY$}
\[
XY = X(-2X+3)= -2X^2 + 3X
\]

\noindent \textbf{Step 2: Compute $E[Y]$}
\[
E[Y]=E[-2X+3]=-2E[X]+3
\]

\noindent \textbf{Step 3: Substitute}
\begin{align*}
\text{Cov}(X,Y) &= E[-2X^2+3X] - E[X](-2E[X]+3) \\
&= -2E[X^2] + 3E[X] + 2(E[X])^2 -3E[X] \\
&= -2E[X^2] + 2(E[X])^2
\end{align*}

\noindent \textbf{Step 4: Recognize variance}
\[
\text{Var}(X)=E[X^2]-(E[X])^2
\]

\noindent Thus:
\[
\text{Cov}(X,Y) = -2\text{Var}(X)
\]

\noindent \textbf{Final Answer:}
\[
\boxed{\text{Cov}(X,Y) = -2\text{Var}(X)}
\]

\newpage

\subsection*{5. Worked Example 3 (Continuous Case from PDF)}

\noindent\textbf{Given (Page 33):}
\[
f_{XY}(x,y)=\frac{xy}{96}, \quad 0<x<4,\;1<y<5
\]

\noindent \textbf{Step 1: Compute $E[X]$}
\[
E[X]=\frac{1}{8}\int_0^4 x^2 dx = \frac{1}{8}\cdot \frac{4^3}{3} = \frac{8}{3}
\]

\noindent \textbf{Step 2: Compute $E[Y]$}
\[
E[Y]=\frac{1}{12}\int_1^5 y^2 dy = \frac{31}{9}
\]

\noindent \textbf{Step 3: Compute $E[XY]$}\\
Since independent:
\[
E[XY]=E[X]E[Y]=\frac{8}{3}\cdot \frac{31}{9}=\frac{248}{27}
\]

\noindent \textbf{Step 4: Compute $E[2X+3Y]$}\\
Using linearity:
\[
E[2X+3Y] = 2E[X]+3E[Y]=\frac{47}{3}
\]

\noindent \textbf{Step 5: Compute Variances}
\[
\text{Var}(X)=\frac{8}{9}, \quad \text{Var}(Y)=\frac{92}{81}
\]

\noindent \textbf{Step 6: Compute Covariance}
\[
\text{Cov}(X,Y)=E[XY]-E[X]E[Y]=0
\]

\noindent \textbf{Step 7: Pearson Coefficient}
\[
\rho_{XY}=0
\]

\newpage

\section*{Topic 4: Pearson’s Correlation Coefficient}

\subsection*{1. Definition}
\[
\rho_{XY} = \frac{\text{Cov}(X,Y)}{\sqrt{\text{Var}(X)\text{Var}(Y)}}
\]

\subsection*{2. Properties (Page 31–32)}
\begin{itemize}
    \item Measures strength + direction
    \item Normalized covariance
    \item Unit-free
    \item Range: $-1 \le \rho \le 1$
\end{itemize}

\noindent\textbf{Interpretation}
\begin{itemize}
    \item (+1): Perfect positive
    \item (-1): Perfect negative
    \item (0): No linear correlation
\end{itemize}

\noindent\textbf{Practice Example (Mitzenmacher \& Upfal Inspired)}\\
If:
\[
\text{Cov}(X,Y)=6,\; \sigma_X=2,\; \sigma_Y=3
\]
\[
\rho=\frac{6}{2\cdot 3}=1
\]
\textbf{Interpretation:} Perfect linear relationship.

\newpage

\section*{Topic 5: Joint Gaussian Random Variables}

\subsection*{1. Definition}
A pair $(X,Y)$ is jointly Gaussian if:
\begin{itemize}
    \item Joint PDF is Gaussian
    \item Marginals are Gaussian
\end{itemize}

\subsection*{2. Algebraic Form}
\[
f_{XY}(x,y)=\frac{1}{2\pi\sigma_X\sigma_Y\sqrt{1-\rho^2}}
\exp\left(
-\frac{1}{2(1-\rho^2)}\left[
\left(\frac{x-\mu_X}{\sigma_X}\right)^2
+
\left(\frac{y-\mu_Y}{\sigma_Y}\right)^2
-2\rho\left(\frac{x-\mu_X}{\sigma_X}\right)\left(\frac{y-\mu_Y}{\sigma_Y}\right)
\right]
\right)
\]

\noindent\textbf{Parameters}
\begin{itemize}
    \item $(\mu_X, \mu_Y)$: Means
    \item $(\sigma_X, \sigma_Y)$: Standard deviations
    \item $(\rho)$: Correlation coefficient
\end{itemize}

\subsection*{3. Interpretation (Page 53–54)}
\noindent\textbf{Example:}\\
$X$: CO₂ concentration\\
$Y$: Temperature

\noindent If:
\[
\rho > 0 \Rightarrow \text{Higher CO₂ → Higher temperature}
\]

\noindent\textbf{Applications:}
\begin{itemize}
    \item Signal processing
    \item Communication systems
    \item Noisy measurements
\end{itemize}

\newpage

\section*{FINAL SUMMARY (MEMORY-OPTIMIZED)}

\subsection*{1. Key Concepts (Flashcard Style)}
\begin{itemize}
    \item \textbf{Marginal PMF} $\rightarrow$ “Sum out the other variable”
    \item \textbf{Correlation} $\rightarrow$ “Raw product measure”
    \item \textbf{Covariance} $\rightarrow$ “Centered dependence”
    \item \textbf{Pearson} $\rightarrow$ “Normalized covariance”
    \item \textbf{Gaussian} $\rightarrow$ “Bell-shaped joint behavior”
\end{itemize}

\subsection*{2. Formulas to Remember}
\begin{align*}
p_X(x) &= \sum_y p_{XY}(x,y) \\
\text{Cov}(X,Y) &= E[XY]-E[X]E[Y] \\
\rho &= \frac{\text{Cov}(X,Y)}{\sigma_X\sigma_Y}
\end{align*}

\subsection*{3. What to Learn}
\begin{itemize}
    \item Difference: correlation vs covariance
    \item Why centering is important
    \item Independence vs zero covariance
    \item Reading joint distributions
\end{itemize}

\subsection*{4. Applications (Exam-Oriented)}
\begin{itemize}
    \item \textbf{Marginalization} $\rightarrow$ tables \& joint distributions
    \item \textbf{Covariance} $\rightarrow$ detect direction of dependence
    \item \textbf{Pearson} $\rightarrow$ measure strength
    \item \textbf{Gaussian} $\rightarrow$ modeling real-world noise
\end{itemize}

\subsection*{5. Tricks}
\begin{itemize}
    \item \textbf{Independence} $\Rightarrow$ directly use $E[XY]=E[X]E[Y]$
    \item \textbf{Linear transformation} $\rightarrow$ use covariance property
    \item \textbf{For tables} $\rightarrow$ sum rows/columns quickly
\end{itemize}

\subsection*{6. Be Careful:}
\begin{itemize}
    \item Zero covariance $\neq$ independence
    \item Correlation depends on scaling
    \item Forgetting to subtract mean in covariance
    \item Misusing joint vs marginal distributions
    \item Integration limits in continuous cases
\end{itemize}

\begin{center}
\textbf{END OF SCRIBE}
\end{center}

\end{document}